\label{sec:methodology}

We first define the algebraic object differentiated by every route. Let
$V_{h,0}=\operatorname{span}\{\varphi_i\}_{i=1}^{n}$ be the constrained
finite-element space, and let $x\in\mathbb{R}^n$ contain its free scalar
coefficients. Essential data enter through an affine lift
$\bar u_h(\theta)$, where $\theta$ collects fixed load and material
parameters. The corresponding state is
\begin{equation}
  u_h(x;\theta)
  =\bar u_h(\theta)+\sum_{i=1}^{n}x_i\varphi_i.
  \label{eq:constrained-state}
\end{equation}
For element $e$, we write the local state as
$u_e(x;\theta)=A_ex+b_e(\theta)$, where $A_e$ is the fixed restriction and
component map. This notation includes nonhomogeneous prescribed values and
arbitrary free-DOF permutations.

Let $B_{eq}$ map the local state to the strain or gradient argument at
quadrature point $q$. We fix the quadrature rule, weights, local material data
$m_{eq}$, and a deterministic first-match selector $\beta_{eq}(x)$ evaluated
from the local strain and data. The discrete element contribution is
\begin{equation}
  \Phi_e(x;\theta)
  :=\sum_{q=1}^{n_q}w_{eq}
  \mathcal{W}_{\beta_{eq}(x)}
  \bigl(B_{eq}(A_ex+b_e);m_{eq}\bigr),
  \label{eq:element-quadrature-potential}
\end{equation}
and the global free-space potential is
\begin{equation}
  \Pi_h(x;\theta)
  :=\sum_{e\in\mathcal{T}_h}\Phi_e(x;\theta)-f_h^\top x.
  \label{eq:generic-potential}
\end{equation}
Additive terms that depend only on prescribed data may be omitted because they
do not alter the free-space derivatives. The symbol $P_k(L_\ell)$ denotes
degree-$k$ Lagrange elements on mesh level $L_\ell$.

\subsection{Element and constitutive differentiation}
\label{subsec:element-constitutive-equivalence}

Element AD applies first- and second-order differentiation directly to
$\Phi_e$. At a fixed state $x_\star$, write
$\beta_{eq}^\star:=\beta_{eq}(x_\star)$. Constitutive AD differentiates that
selected scalar density with respect to
\begin{equation}
  \varepsilon_{eq}:=B_{eq}(A_ex+b_e),
  \qquad
  \sigma_{eq}:=\partial_\varepsilon\mathcal{W}_{\beta_{eq}^\star},
  \qquad
  C_{eq}:=\partial^2_{\varepsilon\varepsilon}
  \mathcal{W}_{\beta_{eq}^\star}.
  \label{eq:constitutive-objects}
\end{equation}
The branch-stability assumption below makes
$\beta_{eq}(x)=\beta_{eq}^\star$ on a neighborhood of $x_\star$; the selector
is therefore locally constant during branch-interior differentiation.

\begin{proposition}[Fixed-functional branchwise equivalence]
\label{prop:branchwise-route-equivalence}
Fix an element state $x_\star$. Assume that, on a neighborhood of $x_\star$,
(i) both routes use the same scalar densities, quadrature weights, material and
history data; (ii) $u_e(x)=A_ex+b_e$ is affine; (iii) every $B_{eq}$ is fixed;
(iv) each selected density is twice continuously differentiable; (v) the
branch labels and principal-value ordering do not change; and (vi) both routes
use the same tensor and engineering-shear conventions. Then element AD and
constitutive AD give identical element residuals and Hessians in exact
arithmetic:
\begin{align}
  \nabla_x\Phi_e(x)
  &=\sum_q w_{eq}A_e^\top B_{eq}^\top\sigma_{eq},
  \nonumber\\
  \nabla_x^2\Phi_e(x)
  &=\sum_q w_{eq}A_e^\top B_{eq}^\top
    C_{eq}B_{eq}A_e.
  \label{eq:branchwise-route-identity}
\end{align}
Consequently, assembling either set of element objects gives the same global
free-space gradient and Hessian of \eqref{eq:generic-potential}.
\end{proposition}

\begin{proof}
The affine map has zero second derivative. Applying the chain rule first to
$\varepsilon_{eq}(x)$ and then to every summand of
\eqref{eq:element-quadrature-potential} gives
\eqref{eq:branchwise-route-identity}. Summation and the fixed free-space
assembly map preserve the identities.
\end{proof}

Proposition~\ref{prop:branchwise-route-equivalence} is conditional. It does not
apply at a branch switch, under an unresolved principal-value ordering, or when
the compared codes change quadrature or constraints. It also establishes the
derivatives of the implemented scalar potential, not equivalence to an
incremental return-mapping model.

\subsection{Colored sparse recovery}
\label{subsec:colored-recovery}

Let $H:=\nabla_x^2\Pi_h(x)$, let $I_r$ be the free rows owned by rank $r$,
and let $\mathcal{P}_r$ contain the structural entries $(i,j)$ of those rows.
The local element set $\mathcal{E}_r$ contains every element incident on
$I_r$, and $J_r$ contains every free column appearing in those elements.  We
write $H^{(r)}$ for the Hessian of the sum over $\mathcal{E}_r$, restricted to
$J_r$.  Because all element contributions incident on an owned row are present,
$H^{(r)}_{ij}=H_{ij}$ for $(i,j)\in\mathcal{P}_r$.

The rank-local column-interference graph joins columns $j,k\in J_r$ whenever
some owned row $i\in I_r$ satisfies
$(i,j),(i,k)\in\mathcal{P}_r$. A local color class
$\mathcal{C}_{r,c}$ defines the overlap seed
\begin{equation}
  v_{r,c}:=\sum_{j\in\mathcal{C}_{r,c}}e_j,
  \qquad y_{r,c}:=H^{(r)}v_{r,c}.
  \label{eq:colored-seed}
\end{equation}
This construction is the direct Hessian analogue of classical sparse
derivative probing \citep{curtis1974sparsejacobian,coleman1984sparsehessian}.

\begin{proposition}[Owned-row colored recovery]
\label{prop:owned-row-colored-recovery}
For every rank $r$, assume that (i) $\mathcal{P}_r$ contains every active
nonzero $(i,j)$ of each owned row; (ii) $\mathcal{E}_r$ contains every element
contributing to those rows, with complete state, affine-lift, material, and
selected-branch data on its overlap; (iii) $J_r$ contains every column in
$\mathcal{P}_r$, and each owned row contains at most one structural entry from
any $\mathcal{C}_{r,c}$; (iv) $H^{(r)}v_{r,c}$ is evaluated exactly with the
complete overlap seed; and (v) every global row has one PETSc owner. Then, only
for $(i,j)\in\mathcal{P}_r$ with $j\in\mathcal{C}_{r,c}$, assigning
$H_{ij}:=(y_{r,c})_i$ recovers every owned structural entry exactly. The union
of the owned rows is the canonical distributed matrix $H$.
\end{proposition}

\begin{proof}
For an owned row $i$, the identity
$(y_{r,c})_i=\sum_{j\in\mathcal{C}_{r,c}}H^{(r)}_{ij}$ contains at most one
structural nonzero by assumption (iii). Assumptions (i), (ii), and (iv) make
the local sum complete and give $H^{(r)}_{ij}=H_{ij}$ on the owned pattern.
The restricted assignment therefore recovers that entry, and unique ownership
assembles each recovered row once.
\end{proof}

The proposition permits different rank-local color labels and counts because
each rank evaluates its own complete incident-element Hessian and inserts only
its owned pattern. It does not permit a missing incident element or structural
edge, incomplete state or seed overlap, inconsistent branch data, or two
same-color columns that meet an owned row. Our production recovery uses AD
Hessian-vector products; a finite-difference gradient probe would add its
truncation error.

\subsection{Cost components and comparison variables}
\label{subsec:cost-model}

Let $m_e$ be the active element DOF count, $s$ the constitutive dimension,
$n_q$ the quadrature-point count, $n_{e,r}^{\mathrm{loc}}$ the overlap-element
count on rank $r$, and $c_r$ its local color count. Because the measured
response is an MPI collective maximum, the structural work proxy must describe
the busiest rank rather than the global number of uniquely owned elements. We
therefore define
\begin{align}
  \chi_{\mathrm{elem}}
  &:=\max_r n_{e,r}^{\mathrm{loc}}n_qm_e^2,
  \nonumber\\
  \chi_{\mathrm{color}}
  &:=\max_r\bigl(n_{e,r}^{\mathrm{loc}}c_r\bigr)n_qm_e,
  \nonumber\\
  \chi_{\mathrm{const}}
  &:=\max_r n_{e,r}^{\mathrm{loc}}n_q
    \bigl(s^2+s^2m_e+sm_e^2\bigr).
  \label{eq:route-work-proxies}
\end{align}
The first two expressions count the dense element-Hessian contribution shape
and the vector work propagated by the exact AD-HVP seeds, respectively. In the
constitutive expression, $s^2$ counts the dense tangent $C_q$, while $s^2m_e$
and $sm_e^2$ count the two contractions used to form $B_q^\top C_qB_q$.
The proxies in \eqref{eq:route-work-proxies} are structural operation-shape
counts, not exact FLOP, cache, or compiler-cost models. We therefore also
decompose the measured route cost as
\begin{equation}
  T
  =T_{\mathrm{differentiate}}+T_{\mathrm{contract}}
   +T_{\mathrm{insert}}+T_{\mathrm{communicate}}
   +T_{\mathrm{cache}}+T_{\mathrm{imbalance}},
  \label{eq:route-cost-decomposition}
\end{equation}
and record matrix nonzeros, overlap size, memory, branch fractions, rank count,
and the applicable $\chi$ from \eqref{eq:route-work-proxies}. A separate
one-factor-at-a-time diagnostic
varies element DOFs, quadrature points, constitutive dimension, color count,
nonzeros per row, message bytes, and imbalance. For each of four stages, an OLS
model fits the log warm median on ranks 1 and 8 using an intercept, log rank,
and the seven log factor ratios; rank 32 is used once for validation. Three
independent allocation blocks are required per rank. Total synthetic work is
held fixed while multi-rank imbalance is varied; nonunit imbalance is
inapplicable at one rank and those two rows are excluded from the fit. The
diagnostic must have median and 90th-percentile absolute percentage errors at
most $50\%$ and $100\%$, respectively. Because its kernels are synthetic rather
than route-faithful, it is used only for descriptive mechanism checks and is
not a production-model feature.

The final response is the logarithm of the collective-max warm-median
tangent-construction time. Its 13 features are three route indicators, three
route-specific cluster-architecture shifts, $\log(1+\chi)$,
$\log(1+n_{\mathrm{nz}}^{\max})$,
$\log(1+n_{\mathrm{ov}}^{\max})$, log rank, plastic fraction, and two
untransformed max-to-mean imbalance ratios.
More precisely, let $\mathcal{A}:=\{\mathrm{elem},\mathrm{color},
\mathrm{const}\}$, let $I_a$ indicate route $a$, and let $K$ indicate the
cluster architecture.
Furthermore, let $n_{\mathrm{nz}}^{\max}$ be the maximum number of owned
matrix nonzeros over ranks, $n_{\mathrm{ov}}^{\max}$ the maximum overlap-DOF
count, $R$ the MPI
rank count, $p_{\mathrm{pl}}$ the plastic quadrature fraction, and
$\iota_e,\iota_o$ the element and overlap max-to-mean ratios. The frozen model
is
\begin{align}
  \log T
  ={}&\sum_{a\in\mathcal{A}}\alpha_a I_a
    +\sum_{a\in\mathcal{A}}\kappa_a K I_a
    +\beta_\chi\log(1+\chi_a)
    +\beta_{\mathrm{nz}}\log(1+n_{\mathrm{nz}}^{\max})
  \nonumber\\
   &+\beta_{\mathrm{ov}}\log(1+n_{\mathrm{ov}}^{\max})
    +\beta_R\log R+\beta_p p_{\mathrm{pl}}
    +\beta_e\iota_e+\beta_o\iota_o.
  \label{eq:route-time-model}
\end{align}
The production OLS fit has no intercept: the three mutually exclusive route
indicators provide the route-specific levels without intercept collinearity.
Each active route--configuration--state--rank slot contributes one equally
weighted row after its independent block medians are aggregated; censored or
incomplete active-route groups are excluded and never imputed. Peak RSS and
tracked allocation are outcomes, not predictors. Local single-node runs and
cluster ranks 1 and 8 form the training set; cluster rank 32 is the untouched
production holdout. Fitting requires at least 48 training rows, 20 holdout rows,
6/4 training/holdout rows per route, and 12/8 distinct training/holdout groups,
with both architectures represented in training. The 13-column design must
have full column rank and condition number at most $10^{10}$. Percentage errors
are computed in seconds after exponentiating the log-time predictions. The
selector requires holdout median and 90th-percentile absolute percentage errors
at most $25\%$ and $50\%$, at least $90\%$ correct
ordering outside a $10\%$ tie band, at least four resolved groups, and at least
two distinct observed winners. A predictive selector is admissible only after
these and all derivative, state, and endpoint gates pass. Coefficient,
prediction, and route-order intervals use 10,000 paired allocation-block
bootstrap resamples at 95\% confidence with the prespecified seed 20260710.
Section~\ref{sec:performance-solver-behavior} shows that the current data do not
meet this condition.

\subsection{Mesh-aware nonlinear accuracy}
\label{subsec:mesh-aware-stopping}

Coefficient norms change with mesh density and therefore do not define a
portable stopping contract. Let $M_h$ be a symmetric positive-definite Riesz
operator on the constrained free space. We use the primal and dual norms
\begin{equation}
  \lVert s\rVert_{M_h}:=(s^\top M_hs)^{1/2},
  \qquad
  \lVert g\rVert_{M_h^{-1}}:=(g^\top M_h^{-1}g)^{1/2}.
  \label{eq:riesz-primal-dual-norms}
\end{equation}
The scalar problems use a row-sum lumped P1 mass operator. Hyperelasticity and
the branch-structured mechanics problem use the tangent of a reference elastic
energy on the identical constrained free space. Before a reference tangent is
admitted, a scale-aware symmetry check and a recorded direct factorization
report zero negative and zero numerical zero pivots under the stated solver
options. This is a numerical inertia check, not an analytic SPD proof. Each
inverse-Riesz solve also passes an independently recomputed true-residual gate.

For nonlinear iteration $k\geq 1$, let
\begin{equation}
  R_k:=\lVert g_k\rVert_{M_h^{-1}},\qquad
  C_k:=\lVert x_k-x_{k-1}\rVert_{M_h},\qquad
  S_k:=\lVert x_k\rVert_{M_h},\qquad
  D_k:=\lvert\mathcal E_h(x_k)-\mathcal E_h(x_{k-1})\rvert.
  \label{eq:absolute-residual-correction}
\end{equation}
We report the dimensionless quantities
\begin{equation}
  r_k^{\mathrm{rel}}:=\frac{R_k}{R_0},
  \qquad
  c_k^{\mathrm{rel}}:=\frac{C_k}{\max(s_0,S_k)},
  \label{eq:relative-residual-correction}
\end{equation}
where $r_k^{\mathrm{rel}}$ is defined only when $R_0>0$. The scale $s_0>0$
has the same units as the primal norm and is either the initial state norm or
an explicit physical scale. With absolute and relative tolerances denoted by
$\tau_R^{\mathrm{abs}}$, $\tau_R^{\mathrm{rel}}$,
$\tau_C^{\mathrm{abs}}$, and $\tau_C^{\mathrm{rel}}$, the maintained scalar,
hyperelastic, and branch-structured Newton drivers accept an updated iterate
only when
\begin{equation}
\begin{split}
  T_k={}&\{x_k\ \text{accepted and all terminal quantities finite}\}
  \\
  &{}\land\{D_k<\tau_E\}
  \land\{C_k<\tau_C^{\mathrm{abs}}
               \ \lor\ c_k^{\mathrm{rel}}<\tau_C^{\mathrm{rel}}\}
  \\
  &{}\land\{R_k<\max(\tau_R^{\mathrm{abs}},
                         \tau_R^{\mathrm{rel}}R_0)\}.
  \label{eq:terminal-boolean-contract}
\end{split}
\end{equation}
Thus the energy-change test is required; it is not a reporting-only
diagnostic. If $R_0=0$ exactly, the initial-relative residual is undefined and
is not used; any finite implementation fallback value is treated as inapplicable.
A zero-step stationary return then requires finite
$\mathcal E_h(x_0)$ and $R_0<\tau_R^{\mathrm{abs}}$; the energy-change and
correction tests are inapplicable because no update has been accepted.

The Boolean logic in \eqref{eq:terminal-boolean-contract} is applied once for
each scalar solve and each hyperelastic load step, while the branch-structured
full solve applies it to its single prescribed load. A multi-step
hyperelastic trajectory is complete only when every requested step is
successful. The branch-structured publication gate additionally recomputes
the endpoint residual and requires its selected-branch and provenance checks;
these admission tests do not replace $T_k$.
For hyperelasticity, the coefficient vector represents the free part of the
deformation map under one fixed affine lift. Its reference-energy norm is used
only within that representation and is not invariant under a change of lift.
The endpoint record contains the selected Riesz operator, scale, free-space
identity, numerical inertia record, inverse-solve diagnostics, and independently
recomputed terminal norms.

The tolerances in the present diagnostic records are case-specific. The
scalar Ginzburg--Landau pair supplies one local tenfold sensitivity check, but
cross-mesh calibration remains open. The reference-elastic hyperelastic and
branch-structured policies have passed operator smokes only; we therefore do
not claim a universal or publication-calibrated stopping threshold.

\subsection{Evidence and timing contract}
\label{subsec:evidence-contract}

Correctness precedes time. A route row first passes the common state hash,
energy, gradient, matrix or multiple action probes, branch margin, Riesz
residual, and correction gates. Timing then uses warmups, balanced route order,
collective-max MPI phases, independent repetitions, and explicit censoring.
Compilation, setup, assembly, preconditioner construction, Krylov solution, and
total time remain separate quantities. A capped or failed solve is not encoded
as a large successful time.

This hierarchy yields four distinct statements: a proved identity under stated
assumptions; verification on finite tested states; a descriptive finite route
map; and an empirical crossover or selector. Evidence for one level does not
imply the next. In particular, near-equal energy does not replace a residual or
tangent comparison, and one tangent probe does not establish complete matrix
equality.
